
% ============================================================================
\clearpage
\section*{v4.9.0-r81 Source-Only: A Joint Action W1, the Matter--Tachyon Pairing, and the Search for a Vacuum Term}
\addcontentsline{toc}{section}{v4.9.0-r81 Source-Only: A Joint Action W1, the Matter--Tachyon Pairing, and the Search for a Vacuum Term}
\label{sec:v490-r81-joint-action-w1}
% ============================================================================

\paragraph{Purpose.}
Layer v4.9.0-r80 closed the comparison-model branch and moved the open question to the level of a joint dynamics
(\S\ref{sec:v490-r80-background-not-sourced}).  On 4 and 5 November 2026 the author took five decisions on such an action, and the
work that followed is recorded here.  The layer contains the first action W1 and its identities, four pre-registered tests (H69,
H70, H71, H72) and one refuted test (H73), a series of explorations on the matter--tachyon pairing, the vacuum term, the Higgs
field as a tachyonic field and the connectors between matter and its tachyon image, and three corrections of earlier statements.
The research notes are \texttt{notes/\allowbreak research/\allowbreak v4\_10\_0\_w1/} (merged as PR~\#65).  Earlier layers are
not rewritten.

\noindent\textbf{Source-only status.}
This v4.9.0-r81 layer is source-only.  It does not build a PDF, does not assign a new DOI,
and does not authorize a public release.  The v4.9.0 Version DOI
\href{https://doi.org/10.5281/zenodo.23534954}{10.5281/zenodo.23534954} and the Concept DOI
\href{https://doi.org/10.5281/zenodo.19210728}{10.5281/zenodo.19210728} are unchanged by it.

\subsection*{v4.9.0-r81 decisions and the action W1}
\label{subsec:v490-r81-action}

\paragraph{Author decisions (4--5 November 2026).}
\begin{enumerate}[leftmargin=2em]
\item The coframe is an independent field; whether solutions satisfy \(E\approx C[\Psi]\) is to be tested later, not assumed.
\item A real action with Dirac adjoint and symmetric kinetic term in frame M is the main line; the complex reading K (stationarity
of \(\mathrm{Re}\,S\) and \(\mathrm{Im}\,S\)) is computed alongside.
\item Einstein--Cartan gravity without a cosmological term, one system for now.
\item The Dixon carrier \(\mathbb C\otimes\HH\otimes\OO\) with the \(U(1)\) generated by \(J=iL_{e_7}\), extended later by colour and the weak
\(SU(2)\); bookkeeping in QAM.
\item The closure projector \(\hat P_{C_7}\) is left open.
\end{enumerate}

\paragraph{Theoremlet v4.9.0-r81.1 (the action W1).}
Exploration.  W1 is \(S=\int e\,[\tfrac1{2\kappa}R(e,\omega)+\tfrac12(\Psi^\dagger i\bar\sigma^\mu D_\mu\Psi+\mathrm{c.c.})-\tfrac14F^2]\) on the sixteen
left-handed Weyl states of one generation with \(\nu_R\) (\S\ref{sec:v490-r42-mm-compensator-chirality}), Lorentz acting by left
multiplication on \(\HH\).
\begin{enumerate}[leftmargin=2em]
\item[(a)] \(U(1)_J\) is anomaly-free on the full carrier: \(\operatorname{Tr}q=\operatorname{Tr}q^3=0\) for \(J\), \(iR_{e_7}\) and \(B-L\).  On the
image of \(C_{\rm ch}\) alone \(B-L\) would be anomalous (\(\operatorname{Tr}q^3=-1.78\)).
\item[(b)] Eliminating the torsion gives \(\mathcal L_{\rm eff}=-\tfrac3{16}\kappa\,j\cdot j\) with the total current of all sixteen states; the
control reproduces the Hehl--Datta value for a Dirac field.  The torsion is blind to the \(U(1)\) charge.
\item[(c)] W1 produces no vacuum term: without a scalar potential no term proportional to \(g_{\mu\nu}\) with constant coefficient arises.
\end{enumerate}

\subsection*{v4.9.0-r81 the complex reading and the matter--tachyon pairing}
\label{subsec:v490-r81-pairing}

\paragraph{Theoremlet v4.9.0-r81.2 (complex reading).}
Exploration, pointwise identities checked to \(2\times10^{-14}\).  With the symmetric kinetic term \(\mathrm{Im}\,S\equiv0\): K adds nothing.  Without
symmetrization \(\mathrm{Im}\,L=\tfrac12D_\mu j^\mu\); with an independent connection, stationarity of \(\mathrm{Im}\,S\) forces \(j^a=0\), hence \(\Psi=0\)
in a single frame, since \(j^0=|\psi|^2\ge0\).  On the image of \(C_{\rm ch}\) the inner and outer units agree up to sign.

\paragraph{Theoremlet v4.9.0-r81.3 (the matter--tachyon gap).}
Exploration.  With the tachyon frame of Convention~F as the continuation \(\eta\to\eta+i\pi/2\) along an axis \(n\), the vector map is
\(\Lambda(i\pi)=-R_n(\pi)\), so a tachyon component carries the current \(j_T=\Lambda(i\pi)j_M\) of negative density.  Then \(j_M+j_T=2j_\perp\):
K only requires that matter flows along the axis, and the remaining equations from \(\mathrm{Im}\,S\) hold.  For frames offset by \(\delta\)
(author: frames need not be synchronous) the transverse current is antiperiodic with period \(2\delta\).

\paragraph{Theoremlet v4.9.0-r81.4 (H69 confirmed).}
Pre-registered (commit \texttt{dbcc40d}).  Flat space, \(A=0\): two counter-propagating left-handed waves plus their tachyon image offset by
\(\delta\), with kinetic sign \(s_T=-1\) from F3.  The current balance holds exactly iff \(\Delta\omega\,\delta\equiv\pi\pmod{2\pi}\) (residual
\(5\times10^{-15}\), violated at \(0\) and \(\pi/2\)), and the total energy--momentum vanishes including cross terms (\(10^{-15}\)).  The pair
carries neither net current nor net energy; the offset fixes a beat frequency.

\paragraph{Theoremlet v4.9.0-r81.5 (H70: the offset must be comoving).}
Pre-registered (commit \texttt{54c354a}, not blind).  On flat FRW with massless Weyl fields, of the couplings \(\delta=c/H\), a fixed conformal
offset, and \(\delta=t-\sqrt{t^2-\bsym^2}\) from Theorem~\ref{thm:expansion}, only the fixed conformal offset (\(\delta\propto a\) in cosmic time) keeps
the pair consistent (H70a confirmed).  The other two violate the Weyl equation of the tachyon image; the third is undefined for
\(t\le\bsym\).  H70b is refuted: wrong couplings leave a spatially homogeneous residual, which is not \(\Lambda\)-like (only
\(\tau\tau\) and \(\tau z\) components) and comes from the violated field equation.

\paragraph{Theoremlet v4.9.0-r81.6 (frame weights and the frame check).}
Explorations.
\begin{enumerate}[leftmargin=2em]
\item[(a)] The neutral pairing has equal weights, the photon-frame mixture P.  With matter-frame weights a homogeneous, traceless,
radiation-like residual \(\propto p_+-p_-=1/\gamma\) remains.  Massless fields never give a \(\Lambda\)-like term.
\item[(b)] On the doublet, the tachyon map \(R=\sigma_y\) and time reversal \(T=\sigma_z\) give the same image on the compensator axis, where
the matter frame lies; both exchange the port weights.  They differ off the axis and on the carrier, where \(T\) conjugates the charges
and \(R\) fixes the sectors (\S\ref{sec:v490-r51-e7-frames}).
\end{enumerate}

\subsection*{v4.9.0-r81 mass, tachyonic modes and the vacuum term}
\label{subsec:v490-r81-vacuum}

\paragraph{Theoremlet v4.9.0-r81.7 (H71: the massive partner).}
Pre-registered (commit \texttt{4a3b802}; (a), (b) not blind).  For \(L=\psi^\dagger i\bar\sigma\partial\psi+s\,\chi^\dagger i\sigma\partial\chi-m(\psi^\dagger\chi+\mathrm{h.c.})\):
(a) \(s=+1\) gives \(\omega^2=k^2+m^2\), \(s=-1\) gives the tachyonic \(\omega^2=k^2-m^2\) with a spacelike current; (b) for \(s=+1\) the
balance fails; (c) for \(s=-1\) it fails as well (refuted).  An antiparticle and a mass-coupled mirror partner are classically the same
Dirac partner; an uncoupled mirror sector (the author's ``mirror particle'') adds positive, dark-matter-like energy.

\paragraph{Theoremlet v4.9.0-r81.8 (H72 and H73: a vacuum term from tachyonic modes).}
For a plane mode \(T_{\mu\nu}=N\,p_\mu p_\nu\); an isotropic incoherent gas has \(w=\sum N k^2/(3\sum N\omega^2)\).
\begin{enumerate}[leftmargin=2em]
\item[(a)] H72 (commit \texttt{bc85519}): a \(\Lambda\)-like term \(T=\rho\eta\) arises from tachyonic modes of both norm signs on at least two
momentum shells (residual \(3\times10^{-16}\)); its height and sign are set by a free weight ratio; frame weights on one shell give none.
\item[(b)] H73 (commit \texttt{623a080}) refuted: the reference states (matter at rest, tachyon at threshold) with matter-frame weights
give \(w=-\bsym^2/3=-0.10\).
\item[(c)] With modes and weights at a common rapidity the mixture is exactly pressureless dust in frame M (\(\rho=m/\cosh\eta\)) and pure
tension with \(\rho=0\) in frame T, for every \(\eta\).  If gravity sees \(a\,T_M+b\,T_T\), then \(w=-b/(3a)\) and
\(\Omega_\Lambda=b/(3a)\): frame M alone gives no dark energy, equal weights give \(\Omega_m=\tfrac23\), \(\Omega_\Lambda=\tfrac13\), the
exchange of the observed fractions, and frame weights give \(-\bsym^2/3\) again.  None of six justified weightings gives the observed
\(\Omega_\Lambda\approx0.685\); ``two of three sectors'' gives \(0.667\) but was proposed after seeing the target and counts as a fit.
\end{enumerate}
\emph{Reading.}  The structure supplies the mechanism of a vacuum term (tension without energy from the tachyonic part), not its size.

\subsection*{v4.9.0-r81 the Higgs field as a tachyonic field and the connectors}
\label{subsec:v490-r81-connectors}

\paragraph{Theoremlet v4.9.0-r81.9 (connectors between matter and its tachyon image).}
Explorations.
\begin{enumerate}[leftmargin=2em]
\item[(a)] The only gauge-invariant mass term of the carrier, the \(\nu^c\) Majorana connector (r44.1(d)), is never tachyonic:
\(\omega^2=k^2+|M|^2\) for both kinetic signs and all phases.
\item[(b)] Sources: the Higgs field is read as tachyonic at the origin (r06.4(b)) and as an observer-frame image
(Remark~\ref{rem:structural-higgs-role}); the connectors with Higgs quantum numbers are r44.1(b).
\item[(c)] For \(p_M=m(\cosh\eta,\sinh\eta)\) and the tachyon image \(p_T=m(\sinh\eta,\pm\cosh\eta)\) one has \(p_M\cdot p_T=0\) for the
co-moving image, for every \(\eta\).  The connector momenta \(p_M\pm p_T\) give: co-moving, lightlike; counter-moving, the pair connector
\(q^2=+2m^2\sinh2\eta\) and the transition connector \(q^2=-2m^2\sinh2\eta\), tachyonic.  A Higgs-like negative mass squared thus follows
kinematically for counter-moving matter and tachyon image.
\item[(d)] The \(\eta\)-independent \(\Sigma\)-mixtures of the four connectors have \(q^2\in\{0,4,8,16\}\,m^2\), from \(q^2=(n^2-k^2)m^2\);
the bit assignment of the masks only renames them.  No string-like fourth level arises; the frame continuation is cyclic of order four.
\end{enumerate}

\subsection*{v4.9.0-r81 corrections}
\label{subsec:v490-r81-corrections}

\begin{enumerate}[leftmargin=2em]
\item \emph{To r44, Reading 1.}  ``Mass is a coupling between matter frame and tachyon frame'' uses the time reversal \(T\).  By r81.6(b) the
reading is half right: mass terms couple a state to a partner with tachyon-frame weights that is also charge conjugate.  The frame
correction list of r51 (Reading~3) omits r44 and is hereby extended to it.
\item \emph{To the note on r81.8(c).}  The value \(w=-\bsym^2/3\) of H73 is not an artefact: it reappears exactly when the two readings
of the consistent mixture are weighted with the frame weights.
\item \emph{To the draft of W1.}  \(\Psi\) is the whole carrier, not the image of \(C_{\rm ch}\); \(C_{\rm ch}\) is the compensator (the
chirality), not \(\hat P_{C_7}\), which is distinct by the section ``Testing the Native Operator Link'' of the 4.10 research candidate.  The condition \(\mathrm{Tr}(\mathcal M_{16}\hat P_{C_7})\equiv0
\pmod{\mathrm{HC}}\) is not defined as a computation anywhere in the sources.
\end{enumerate}

\paragraph{Proof and measurement sources.}
Scripts \texttt{k170}--\texttt{k187} and the notes in \texttt{notes/\allowbreak research/\allowbreak v4\_10\_0\_w1/}; pre-registrations and
result notes H69--H73 in \texttt{notes/\allowbreak research/\allowbreak predictions/}.  Deferred ideas with triggers (several loop systems
joined by a string, the Chinese remainder reading of ``mod HC'', the dark-energy reading of OP5, condensation of tachyonic modes) are
listed in \texttt{Zurueckgestellte\_\allowbreak Ideen.md} there.  \(\square\)

\paragraph{Consequence.}
A joint action with independent coframe exists in the form W1 and is consistent with the Dixon carrier.  The complex reading forces a
current balance between matter and its tachyon image, which a comoving offset pair satisfies exactly.  A vacuum term needs tachyonic
modes of both norm signs; its size is not fixed by the structure.  OP5 is not closed.

\subsection*{v4.9.0-r81 status in the three frames}
\label{subsec:v490-r81-frames}

\begin{itemize}[leftmargin=2em]
\item W1 is written in frame M.  The tachyon frame T enters only through the declared map \(\eta\to\eta+i\pi/2\) and the sign F3; the
photon frame P is the equal-weight mixture.  All comparisons between M and T are made after this map.
\item r81.6(b) separates the tachyon map \(R\) from time reversal \(T\); results stated for one hold for the other only on the compensator axis.
\item Colour and charge labels are those of the carrier in frame M; they are frame-free between M and T (\S\ref{sec:v490-r51-e7-frames}).
\end{itemize}

\subsection*{v4.9.0-r81 register entries}
\label{subsec:v490-r81-register}

\begin{center}
\small
\begin{tabular}{@{}>{\raggedright\arraybackslash}p{0.34\textwidth}>{\raggedright\arraybackslash}p{0.58\textwidth}@{}}
\toprule
statement & status after v4.9.0-r81\\
\midrule
joint action W1 & set up; anomaly-free \(U(1)_J\); torsion \(-\tfrac3{16}\kappa j^2\); no vacuum term\\
offset matter--tachyon pair & neutral in current and energy (H69); offset comoving (H70)\\
massive tachyon partner & tachyonic dispersion (H71); no current balance with mass\\
vacuum term & possible from tachyonic modes of both norms (H72); size not fixed; H73 refuted\\
Higgs-like tachyonic connector & follows kinematically for counter-moving matter and tachyon image (exploration)\\
r44 reading 1 & corrected (time reversal, not tachyon map)\\
\bottomrule
\end{tabular}
\end{center}

\subsection*{v4.9.0-r81 ledger and guard update}
\label{subsec:v490-r81-ledger-guard}

\[
R_{\rm chk}^{v490\text{-}r81}=(1,0,0)
\]
for (joint action set up and tested in part; vacuum term not fixed; OP5 not closed).

The layer spends no guard:
\[
R_{\rm glob}^{(v490\text{-}r81,+)}=(0,1,1,1,1),
\qquad
\Gamma_{\rm guard}^{v490\text{-}r81}=(1,1,1,1,1),
\qquad
\beta_4=1/256\ \text{unspent}.
\]
It does not imply exact finite \(C5\) zero closure, exact all-readout terminality,
\(C6/T2\), a global photon-clock/RH theorem, or global all-readout closure.
